% GENERATED FILE. Edit canonical lessons, then run npm run generate:books.
% canonical-source-hash: fnv1a64:0cfd418a464465f5
% canonical-lessons: ZH-C502-youhua, ZH-C502-guohua, ZH-W502-hua-paint, ZH-C502-huajia, ZH-C502-minghua, ZH-C502-huazhan

\chapter{Painting}
\label{ch:zh-painting}
\hlchaptermodality{502}

\begin{chapteropening}
\textbf{By the end of this chapter:} \emph{I can talk about oil painting, Chinese painting, painters, famous paintings and art shows.}

Complete the last lesson of chapter 502: I can talk about oil painting, Chinese painting, painters, famous paintings and art shows.
\end{chapteropening}

\section[yóuhuà]{yóuhuà — oil painting}

\label{lesson:ZH-C502-youhua}



\begin{quote}
Before the new one: say the Mandarin for to win (a prize), then the Mandarin for a star player.
\end{quote}

\subsection*{You'll want to know: yóuhuà}
\emph{yóuhuà} — \textquotedblleft{}oil painting\textquotedblright{}.

\begin{quote}
\emph{wǒ xué le liǎng nián yóuhuà} — I studied oil painting for two years.
\end{quote}

Its characters come later; for now it is a word for the ear.

\begin{cousinweb}
\emph{yóuhuà} is built from \emph{yóu} 'oil' and \emph{huà} 'a painting', both inherited from Middle Chinese, \emph{yuw} and \emph{hweaH} (Guangyun, in Baxter's transcription), which glosses the second, after an older book, 'to hang': to hang up the shapes of things in the five colours.
\end{cousinweb}

\begin{grammarlens}[title={What you've built}]
One more word for saying how it was.
\end{grammarlens}

\subsection*{Your turn}
\begin{itemize}
\raggedright
  \item \emph{Say it:} \emph{yóuhuà}
  \item \emph{Say it:} \emph{yóuhuà}, once more
  \item \emph{Say it:} \emph{wǒ xué le liǎng nián yóuhuà} — I studied oil painting for two years.
  \item \emph{Recall:} say the Mandarin for to win (a prize), then the Mandarin for a star player, then say \emph{yóuhuà} again
\end{itemize}

\begin{tcolorbox}[breakable,colback=teal!4,colframe=teal!35!black,title={Before you move on}]
What does yóuhuà mean? (\textquotedblleft{}Oil painting\textquotedblright{}.) Say it once more.
\end{tcolorbox}
\section[\texorpdfstring{\zh{国画} (guóhuà)}{国画 (guóhuà)}]{\zh{国画} (guóhuà) — traditional Chinese painting}

\label{lesson:ZH-C502-guohua}



\begin{quote}
Before the new one: say the Mandarin for a star player, then the Mandarin for an oil painting.
\end{quote}

\subsection*{You'll want to know: \zh{国画}}
\textbf{\zh{国画}} — \emph{guóhuà} — \textquotedblleft{}traditional Chinese painting\textquotedblright{}.

\begin{quote}
\emph{wǒ juéde guóhuà bǐ yóuhuà yǒuyìsi} — I find Chinese painting more interesting than oil painting.
\end{quote}

The new character, in pieces: \zh{画} — eight strokes: a bar on top, then five strokes as on top of \zh{男} \emph{nán}, then an open box of two strokes around them. The next lesson puts it in your hand.

\begin{cousinweb}
\textbf{\zh{国画}} (\emph{guóhuà}) is built from \zh{国} 'a country' and \zh{画} 'a painting', both inherited from Middle Chinese, \emph{kwok} and \emph{hweaH} (Guangyun, in Baxter's transcription), which glosses the first 'a state'.
\end{cousinweb}

\begin{grammarlens}[title={What you've built}]
One more word for saying how it was.
\end{grammarlens}

\subsection*{Your turn}
\begin{itemize}
\raggedright
  \item \emph{Say it:} \emph{guóhuà}
  \item \emph{Say it:} \emph{guóhuà}, once more
  \item \emph{Say it:} \emph{wǒ juéde guóhuà bǐ yóuhuà yǒuyìsi} — I find Chinese painting more interesting than oil painting.
  \item \emph{Recall:} say the Mandarin for a star player, then the Mandarin for an oil painting, then say \emph{guóhuà} again
\end{itemize}

\begin{tcolorbox}[breakable,colback=teal!4,colframe=teal!35!black,title={Before you move on}]
What does \zh{国画} mean? (\textquotedblleft{}Traditional Chinese painting\textquotedblright{}.) Say it once more.
\end{tcolorbox}
\section[\texorpdfstring{\zh{画} (huà)}{画 (huà)}]{\zh{画} — a bar, a field, an open box}

\label{lesson:ZH-W502-hua-paint}



\begin{quote}
The word first, then the character inside it.

\begin{itemize}
\raggedright
  \item \emph{Recall:} say \emph{guóhuà} — \textbf{R1}, one lesson back, before it decays
\end{itemize}

\begin{quote}
\zh{男} \zh{国}
\end{quote}

\emph{Read it:} each one aloud, with nothing said first — \emph{nán}, \emph{guó}
\end{quote}

\subsection*{Script — a bar, a field, an open box}
\hlblockfigure{\detokenize{figures/ZH-W502-hua-paint-filmstrip.pdf}}{How \zh{画} is written, stroke by stroke}

\begin{quote}
\zh{画}
\end{quote}

\textbf{A bar on top.}

\textbf{Then five strokes, as on top of \zh{男}}: a vertical, a bar that turns down, a middle bar, a middle vertical, and the bottom bar.

\textbf{Then the open box, two strokes}: one that runs down the left and turns right along the bottom, and last a vertical on the right, from top to bottom.

1 + 5 + 2 = eight.

\subsection*{Your turn}
\begin{itemize}
\raggedright
  \item \emph{Write it:} \zh{画}, all eight, with the model in front of you
  \item \emph{Cover:} write it again from memory, and count the pieces, not the strokes
\end{itemize}

\begin{tcolorbox}[breakable,colback=teal!4,colframe=teal!35!black,title={Before you move on}]
How many strokes does \zh{画} have? (\textbf{Eight}.) What sits in the middle? (\textbf{The top of \zh{男}}.) Which stroke is last? (\textbf{The vertical on the right}.)

Source: \href{https://raw.githubusercontent.com/chanind/hanzi-writer-data/68d10a4b21150cae5e1ebbd223eed289cf32d90c/data/\%E7\%94\%BB.json}{Hanzi Writer Data \zh{画}, pinned stroke paths and medians}.
\end{tcolorbox}
\section[\texorpdfstring{\zh{画家} (huàjiā)}{画家 (huàjiā)}]{\zh{画家} (huàjiā) — painter (an artist)}

\label{lesson:ZH-C502-huajia}



\begin{quote}
Before the new one: say the Mandarin for an oil painting, then the Mandarin for a traditional Chinese painting.
\end{quote}

\subsection*{You'll want to know: \zh{画家}}
\textbf{\zh{画家}} — \emph{huàjiā} — \textquotedblleft{}painter (an artist)\textquotedblright{}.

\begin{quote}
\textbf{\zh{小时候我想当画家}} \emph{xiǎoshíhou wǒ xiǎng dāng huàjiā} — When I was little I wanted to be a painter.
\end{quote}

\begin{cousinweb}
\textbf{\zh{画家}} (\emph{huàjiā}) is built from \zh{画} 'a painting' and \zh{家} 'a household', both inherited from Middle Chinese, \emph{hweaH} and \emph{kae} (Guangyun, in Baxter's transcription), which glosses the second 'a dwelling'; the sense 'a master of an art' is not given there.
\end{cousinweb}

\begin{grammarlens}[title={What you've built}]
One more word for saying how it was.
\end{grammarlens}

\subsection*{Your turn}
\begin{itemize}
\raggedright
  \item \emph{Say it:} \emph{huàjiā}
  \item \emph{Say it:} \emph{huàjiā}, once more
  \item \emph{Say it:} \emph{xiǎoshíhou wǒ xiǎng dāng huàjiā} — When I was little I wanted to be a painter.
  \item \emph{Recall:} say the Mandarin for an oil painting, then the Mandarin for a traditional Chinese painting, then say \emph{huàjiā} again
\end{itemize}

\begin{tcolorbox}[breakable,colback=teal!4,colframe=teal!35!black,title={Before you move on}]
What does \zh{画家} mean? (\textquotedblleft{}Painter (an artist)\textquotedblright{}.) Say it once more.
\end{tcolorbox}
\section[\texorpdfstring{\zh{名画} (mínghuà)}{名画 (mínghuà)}]{\zh{名画} (mínghuà) — famous painting}

\label{lesson:ZH-C502-minghua}



\begin{quote}
Before the new one: say the Mandarin for a traditional Chinese painting, then the Mandarin for a painter (an artist).
\end{quote}

\subsection*{You'll want to know: \zh{名画}}
\textbf{\zh{名画}} — \emph{mínghuà} — \textquotedblleft{}famous painting\textquotedblright{}.

\begin{quote}
\textbf{\zh{这些名画我都看过}} \emph{zhèxiē mínghuà wǒ dōu kàn guo} — I've seen all these famous paintings.
\end{quote}

\begin{cousinweb}
\textbf{\zh{名画}} (\emph{mínghuà}) is built from \zh{名} 'a name' and \zh{画} 'a painting', both inherited from Middle Chinese, \emph{mjieng} and \emph{hweaH} (Guangyun, in Baxter's transcription), which glosses the first 'a name' and 'a title'.
\end{cousinweb}

\begin{grammarlens}[title={What you've built}]
One more word for saying how it was.
\end{grammarlens}

\subsection*{Your turn}
\begin{itemize}
\raggedright
  \item \emph{Say it:} \emph{mínghuà}
  \item \emph{Say it:} \emph{mínghuà}, once more
  \item \emph{Say it:} \emph{zhèxiē mínghuà wǒ dōu kàn guo} — I've seen all these famous paintings.
  \item \emph{Recall:} say the Mandarin for a traditional Chinese painting, then the Mandarin for a painter (an artist), then say \emph{mínghuà} again
\end{itemize}

\begin{tcolorbox}[breakable,colback=teal!4,colframe=teal!35!black,title={Before you move on}]
What does \zh{名画} mean? (\textquotedblleft{}Famous painting\textquotedblright{}.) Say it once more.
\end{tcolorbox}
\section[huàzhǎn]{huàzhǎn — art exhibition}

\label{lesson:ZH-C502-huazhan}



\begin{quote}
Before the new one: say the Mandarin for a painter (an artist), then the Mandarin for a famous painting.
\end{quote}

\subsection*{You'll want to know: huàzhǎn}
\emph{huàzhǎn} — \textquotedblleft{}art exhibition\textquotedblright{}.

\begin{quote}
\emph{xīngqīliù wǒmen qù kàn huàzhǎn} — We're going to an art exhibition on Saturday.
\end{quote}

Its characters come later; for now it is a word for the ear.

\begin{cousinweb}
\emph{huàzhǎn} is built from \zh{画} 'a painting' and \emph{zhǎn} 'to unfold', both inherited from Middle Chinese, \emph{hweaH} and \emph{trjenX} (Guangyun, in Baxter's transcription), which glosses the second 'to spread out'.
\end{cousinweb}

\begin{grammarlens}[title={What you've built}]
One more word for saying how it was.
\end{grammarlens}

\subsection*{Your turn}
\begin{itemize}
\raggedright
  \item \emph{Say it:} \emph{huàzhǎn}
  \item \emph{Say it:} \emph{huàzhǎn}, once more
  \item \emph{Say it:} \emph{xīngqīliù wǒmen qù kàn huàzhǎn} — We're going to an art exhibition on Saturday.
  \item \emph{Recall:} say the Mandarin for a painter (an artist), then the Mandarin for a famous painting, then say \emph{huàzhǎn} again
\end{itemize}

\begin{tcolorbox}[breakable,colback=teal!4,colframe=teal!35!black,title={Before you move on}]
What does huàzhǎn mean? (\textquotedblleft{}Art exhibition\textquotedblright{}.) Say it once more.
\end{tcolorbox}
